\documentclass[11pt]{article}
\usepackage[a4paper,margin=0.9in]{geometry}
\usepackage{amsmath,amssymb,amsthm,booktabs,enumitem,microtype}
\usepackage[T1]{fontenc}
\usepackage{lmodern}
\usepackage[hidelinks]{hyperref}

\newtheorem{theorem}{Theorem}
\newtheorem{lemma}{Lemma}
\newtheorem{corollary}{Corollary}
\newtheorem{proposition}{Proposition}

\title{Cutoff-One Shared-Clock Obstruction\\
\large Exact Prime-Row Orthogonality in a Literal Two-Channel Audit}
\author{Liang Wang\\
\small Huazhong University of Science and Technology\\
\small Wuhan, China}
\date{August 22, 2026}

\begin{document}
\maketitle

\begin{abstract}
TPC-224 showed that the prime-AP and polarized channels can be written on one
common literal Hilbert family, but left open whether both marginals can save
energy on one source clock.  This paper audits the named finite clock
\(x=Q^3\), \(H=4Q^2\), \(h=4Q\), with primes \(Q<q\leq 2Q\).  Its literal
cutoff is exactly one.  We prove that the two residue coordinates of distinct
prime rows are disjoint, for arbitrary packet profile values.  Consequently
\(E_{\rm AP}=E_{\rm diag}\) and \(E_{\rm all}=E_{\rm pol}\) exactly.  Thus this
clock cannot pay any positive prime-label AP saving; any cancellation is
entirely packet-direction and profile-dependent.  Nine exact-rational affine
scales, aligned and balanced fixtures, and a complete \(Q=3,\ldots,99\)
boundary replay confirm the identities.  The result is a structural,
scoped obstruction, not an asymptotic arithmetic \(L^2\) theorem.
\end{abstract}

\section{Question and claim ceiling}

The previous compatibility audit \cite{tpc224} established that a common
family \(W_{q,j}\) supports three energies
\begin{equation}
 E_{\rm diag}=\sum_{q,j}\|W_{q,j}\|^2,\qquad
 E_{\rm AP}=\sum_j\left\|\sum_qW_{q,j}\right\|^2,
 \qquad
 E_{\rm pol}=\sum_q\left\|\sum_jW_{q,j}\right\|^2,
 \label{eq:energies}
\end{equation}
and that \(E_{\rm all}\leq PJ(E_{\rm AP}+E_{\rm pol})/(P+J)\).  The next
research question is narrower: does the source clock used for the finite
audit actually create prime-label overlap on which an AP dispersion estimate
could operate?

We answer this question exactly for the declared clock.  Our claim is not
that every possible V46 clock has the same geometry.  The clock is a finite
modeling choice inherited from the preceding audit, and no identification
with the physical fixed atom is made.  The principal result is therefore
classified as \texttt{PROVED\_STRUCTURAL\_L1} with a scoped obstruction.

\section{The literal cutoff-one object}

Let \(Q\geq 3\) be an integer and let
\(\mathcal Q_Q=\{q: q\text{ prime},\ Q<q\leq 2Q\}\).  Set
\begin{equation}
 x=Q^3,\qquad H=4Q^2,\qquad h=4Q,
 \qquad C_h=\frac1h.
 \label{eq:clock}
\end{equation}
For a packet profile \(\psi_j\), define the literal residue vector
\begin{equation}
 W_{q,j}(a)=C_h\sum_{0<|m|\leq\lfloor hq/H\rfloor}
 \psi_j\!\left(\frac{Hm}{hq}\right)
 \mathbf 1_{m q^{-1}\equiv a\pmod h}.
 \label{eq:row}
\end{equation}
The inverse exists because an active prime is odd and larger than every prime
factor of \(Q\).  The vector space is the Euclidean space indexed by residues
modulo \(h\).  All four energies in \eqref{eq:energies}, and
\(E_{\rm all}=\|\sum_{q,j}W_{q,j}\|^2\), use exactly this vector and this
normalization.

The clock has a decisive feature:
\begin{equation}
 \left\lfloor\frac{hq}{H}\right\rfloor
 =\left\lfloor\frac qQ\right\rfloor=1.
 \label{eq:cutoff}
\end{equation}
The endpoint \(q=2Q\) is not prime for \(Q\geq3\), so no exceptional endpoint
is hidden in \eqref{eq:cutoff}.  Writing \(r_q=q^{-1}\pmod{4Q}\), every row
has the form
\begin{equation}
 W_{q,j}=u_{q,j}e_{r_q}+v_{q,j}e_{-r_q},\qquad
 u_{q,j}=C_h\psi_j(Q/q),\quad v_{q,j}=C_h\psi_j(-Q/q).
 \label{eq:twopoint}
\end{equation}

\begin{lemma}[disjoint inverse-residue supports]
For distinct \(q_1,q_2\in\mathcal Q_Q\),
\(\{r_{q_1},-r_{q_1}\}\cap\{r_{q_2},-r_{q_2}\}=\varnothing\).
\end{lemma}
\begin{proof}
An intersection would give
\(q_1^{-1}\equiv\varepsilon q_2^{-1}\pmod{4Q}\) for
\(\varepsilon\in\{1,-1\}\).  Multiplication by the units \(q_1q_2\)
gives \(q_2\equiv\varepsilon q_1\pmod{4Q}\).  For \(\varepsilon=1\),
the absolute difference is smaller than \(Q<4Q\), hence the primes are equal.
For \(\varepsilon=-1\), divisibility by \(4Q\) and
\(2Q<q_1+q_2\leq4Q\) force \(q_1+q_2=4Q\).  The upper bounds then force
both primes to equal \(2Q\), which is not prime.  Both alternatives are
contradictions.
\end{proof}

\section{Exact marginal identities}

\begin{theorem}[cutoff-one shared-clock obstruction]
For every finite collection of real packet profiles on the clock
\eqref{eq:clock}, the energies satisfy
\begin{equation}
 E_{\rm AP}=E_{\rm diag},\qquad E_{\rm all}=E_{\rm pol}.
 \label{eq:main}
\end{equation}
In particular, if \(E_{\rm diag}>0\), then no \(\delta>0\) can satisfy
\(E_{\rm AP}\leq(1-\delta)E_{\rm diag}\) uniformly on this clock.
\end{theorem}
\begin{proof}
The lemma makes the rows from distinct prime labels orthogonal, including
their supports across all packet labels.  For fixed \(j\), Pythagoras gives
\[
 \left\|\sum_{q\in\mathcal Q_Q}W_{q,j}\right\|^2
 =\sum_{q\in\mathcal Q_Q}\|W_{q,j}\|^2.
\]
Summing this equality over \(j\) proves the first identity.  For the second,
let \(Y_q=\sum_jW_{q,j}\).  Each \(Y_q\) is supported on the two-point set
of its prime label, so the \(Y_q\) are pairwise orthogonal.  Therefore
\[
 \left\|\sum_qY_q\right\|^2=\sum_q\|Y_q\|^2,
\]
which is exactly \(E_{\rm all}=E_{\rm pol}\).  The final statement follows
from the first identity when the diagonal energy is nonzero.
\end{proof}

The theorem is stronger than a numerical observation: no choice of finite
profile values can create prime-label cancellation before a nontrivial
cutoff or a different, source-identified transform is introduced.  It also
separates the two channel roles.  The AP marginal is rigidly diagonal, while
the polarized marginal may vary from complete packet cancellation to maximal
packet alignment.

\section{Affine profiles and exact closed forms}

For the affine profiles used in TPC-224, write
\(\psi_j(t)=1+s_jt\), \(t_q=Q/q\), and set
\(S_1=\sum_js_j\), \(S_2=\sum_js_j^2\).  Directly from
\eqref{eq:twopoint},
\begin{align}
 E_{\rm diag}=E_{\rm AP}
  &=2C_h^2\sum_{q\in\mathcal Q_Q}(J+S_2t_q^2),
 \label{eq:affined}\\
 E_{\rm pol}=E_{\rm all}
  &=2C_h^2\sum_{q\in\mathcal Q_Q}(J^2+S_1^2t_q^2).
 \label{eq:affinep}
\end{align}
For \(J=4\) and \(s=(0,1,-1,2)/10\),
\(S_1=1/5\), \(S_2=3/50\), so the summands are
\begin{equation}
 C_h^2\left(8+\frac{3}{25}t_q^2\right),\qquad
 C_h^2\left(32+\frac{2}{25}t_q^2\right).
 \label{eq:special}
\end{equation}
The second summand is near four times the first on the source shell.  This
is packet alignment, not an AP saving.

\begin{table}[t]
\centering
\caption{Exact-rational affine source-clock replay.  The two identity columns
are exactly one at every scale; the final column is shown to six decimals only
for readability and is not used as evidence.}
\label{tab:affine}
\begin{tabular}{rrrrr}
\toprule
\(Q\) & \(\#\mathcal Q_Q\) & \(E_{\rm AP}/E_{\rm diag}\)
 & \(E_{\rm all}/E_{\rm pol}\) & \(E_{\rm pol}/E_{\rm diag}\)\\
\midrule
11  & 3  & 1 & 1 & 3.975681\\
17  & 4  & 1 & 1 & 3.975294\\
29  & 6  & 1 & 1 & 3.974163\\
43  & 9  & 1 & 1 & 3.976754\\
61  & 12 & 1 & 1 & 3.975547\\
89  & 16 & 1 & 1 & 3.975183\\
127 & 23 & 1 & 1 & 3.975056\\
181 & 30 & 1 & 1 & 3.975418\\
257 & 42 & 1 & 1 & 3.975244\\
\bottomrule
\end{tabular}
\end{table}

\section{Adversarial profile controls}

The exact obstruction does not depend on the affine choice.  For an aligned
profile, all four plus and minus values are one.  Each prime row then has
\(E_{\rm pol}/E_{\rm diag}=4\), and \(E_{\rm all}=E_{\rm pol}\).  For a
balanced profile, choose plus values \((1,-1,0,0)\) and minus values
\((0,0,1,-1)\).  The packet sums vanish at every prime, so
\begin{equation}
 E_{\rm pol}=E_{\rm all}=0,\qquad E_{\rm AP}=E_{\rm diag}>0.
 \label{eq:balanced}
\end{equation}
These fixtures bracket the packet direction while leaving the AP identity
unchanged.  A complete boundary replay for \(3\leq Q\leq99\) found no
support collision and verified \eqref{eq:main} with exact integers and
rationals.

\section{Block decomposition and certification protocol}

The support lemma has a useful operator interpretation.  Let
\[
 {\cal R}_q=\operatorname{span}\{e_{r_q},e_{-r_q}\}
 \quad\text{and}\quad
 {\cal H}_Q=\mathbb R^{4Q}.
\]
For the named clock, the active part of the Hilbert space is the orthogonal
direct sum
\begin{equation}
 \bigoplus_{q\in{\cal Q}_Q}{\cal R}_q\ \subseteq\ {\cal H}_Q .
\label{eq:block}
\end{equation}
Thus the source object is not merely a collection of rows with small
support: it is a block-diagonal embedding of the prime label into residue
space.  Any linear operation that sums over \(q\) while retaining the
literal coordinates acts independently on these blocks.  The AP marginal
therefore has no cross-block terms to estimate, and the full sum has no
cross-block terms after packet aggregation either.  This is the reusable
structure extracted from the obstruction.

The finite replay was designed to distinguish an identity from a floating
point coincidence.  The producer constructs all coefficients as rational
numbers, forms the four energies by exact inner products, and serializes a
canonical JSON certificate.  The independent checker has its own prime
enumerator and row constructor; it does not import the producer.  It checks
the cutoff, the residue support, all energy identities, and the profile
fixtures before comparing the certificate.  Running the independent checker
with and without Python optimization produced byte-identical output.

\begin{table}[t]
\centering
\caption{Exact replay layers and their claim ceilings.}
\label{tab:audits}
\begin{tabular}{p{0.20\linewidth}p{0.25\linewidth}p{0.42\linewidth}}
\toprule
Audit layer & Input scales & Checked statement\\
\midrule
Affine source clock &
\(\begin{gathered}11,17,29,43,61,89\\127,181,257\end{gathered}\) &
\(E_{\rm AP}=E_{\rm diag}\), \(E_{\rm all}=E_{\rm pol}\), and the
closed forms in \eqref{eq:affined}--\eqref{eq:affinep}.\\
Boundary profiles &
\(Q=3,5,7,13,31,47,73\) &
Aligned and balanced packet directions retain the two exact identities;
the balanced diagonal energy is nonzero.\\
Boundary geometry &
Every integer \(3\leq Q\leq99\) &
The cutoff is one and all active prime supports are pairwise disjoint.\\
\bottomrule
\end{tabular}
\end{table}

The replay has a deliberately narrow interpretation.  It certifies the
literal finite object and its exact theorem; it does not estimate how often
primes occur in a longer interval, and it does not replace a source-lock
argument for a physical V46 transform.  In particular, changing \(H/h\) so
that a second \(m\)-layer enters changes the collision graph and invalidates
the block decomposition \eqref{eq:block} as a proof shortcut.  That is why
the next candidate is a nontrivial-cutoff audit rather than an immediate
claim of arithmetic dispersion.

\section{Route consequence}

The shared-clock question from TPC-224 therefore splits into two logically
different tasks.  On the cutoff-one clock, the AP task is stopped: the literal
prime rows are already orthogonal, so a prime-label dispersion saving cannot
come from summing these rows.  The polarized task remains profile-dependent,
and the exact V46 transfer is still open.  A productive AP route must either
use a source-locked clock with \(\lfloor hq/H\rfloor\geq2\), where distinct
prime rows can have genuine multiplicative collisions, or prove a different
physical synthesis map that creates legitimate overlap.  Neither option is
asserted here.

\begin{verbatim}
TPC225_CUTOFF_ONE = PROVED_EXACT
TPC225_SUPPORT_DISJOINTNESS = PROVED_EXACT
TPC225_AP_EQUALS_DIAGONAL = PROVED_EXACT
TPC225_ALL_EQUALS_POLARIZED = PROVED_EXACT
E_AP = E_diag
E_all = E_pol
TPC225_AP_SAVING_ON_NAMED_CLOCK = REFUTED_SCOPED
TPC225_ARITHMETIC_ADVANCE = NO
TPC225_L2 = NONE
TPC225_FULL_GATE_B = OPEN
\end{verbatim}

\section{Conclusion}

TPC-225 supplies a concrete obstruction immediately downstream of the common
Hilbert interface.  In the named cutoff-one source clock, the AP marginal is
exactly the diagonal energy and the full energy is exactly the polarized
marginal.  Thus the next bridge cannot be crossed by simply declaring a
shared clock; it must expose nontrivial cutoff support or a separately proved
source-identified overlap mechanism.  The result is structural L1 only:
there is no arithmetic \(L^2\) advance, fixed-atom credit, strict
\(1/400\) payment, or twin-prime conclusion.

\bibliographystyle{plain}
\bibliography{references}
\end{document}
